%%===============================================================================
%%      Beamer class source. (original from:  Sascha Frank)
%%
%%             File:
%%           Author:    Christoph L. Schneider
%%             Date:    Apr 2008
%%          Compile:    DVItoPS, PStoPDF
%%===============================================================================

\documentclass{beamer}	%use beamer class

\usepackage{german}
\usepackage{graphicx}

\usepackage[absolute,overlay]{textpos}
\setlength{\TPHorizModule}{\paperwidth}
\setlength{\TPVertModule}{\paperheight}


%%%%%%%%%%%%%%%%%%%%%%%%%%%%%%%%%%%%%%%%%%%%%%%%%%%%%%%%%
% use style file 'Leipzig' for colors, fonts and layout %
% leave those options empty, you do NOT wish		%
							%
% OPTIONS						%
% (a) turn on page numbers at the bottom line: numbers	%
% (b) turn on flowchart bar at the top line: flowchart	%
% (c) turn on navigation bar at the bottom line: nav	%
							%
\usetheme[numbers,flowchart]{Leipzig}			%
							%
% e.g. set all options					%
%\usetheme[numbers,flowchart,nav]{Leipzig}		%
% or only numbers					%
%\usetheme[numbers]{Leipzig}				%
%%%%%%%%%%%%%%%%%%%%%%%%%%%%%%%%%%%%%%%%%%%%%%%%%%%%%%%%



\begin{document}

\title{Fancy titles are always cool}   
\author{Christoph Schneider}
\institute{UFZ}
\date{\today}

%this graphic appears on the title slide
%\titlegraphic{\includegraphics[width=2cm]{ufz_logo}}

%this graphic is your background
\usebackgroundtemplate{\includegraphics[width=\paperwidth,height=\paperheight]{ufz_background}}

\begin{frame}
\titlepage
\end{frame}

\begin{frame}
\frametitle{Table of contents}\tableofcontents
\end{frame}


\section{Section Nr.1} 
\begin{frame}\frametitle{Title} 
Every single frame should include a title.
\end{frame}

\subsection{Subsection Nr.1.1}
\begin{frame}\frametitle{Test title}
...cause without titles something is missing.
\end{frame}


\section{Section Nr.2} 
\subsection{List I}
\begin{frame}\frametitle{Itemization}
\begin{itemize}
\item Introduction course to \LaTeX  
\item Advanced course
\item Seminar work and presentations by \LaTeX 
\item Beamerclass
\end{itemize} 
\end{frame}

\begin{frame}\frametitle{Itemization with pauses}
\begin{itemize}
\item Introduction course to \LaTeX  \pause 
\item Advanced course \pause 
\item Seminar work and presentations by \LaTeX  \pause 
\item Beamerclass
\end{itemize} 
\end{frame}

\subsection{List II}
\begin{frame}\frametitle{Numerated List}
\begin{enumerate}
\item Introduction course to \LaTeX  
\item Advanced course
\item Seminar work and presentations by \LaTeX 
\item Beamerclass
\end{enumerate}
\end{frame}

\begin{frame}\frametitle{Numerated Lists with pauses}
\begin{enumerate}
\item Introduction course to \LaTeX  \pause 
\item Advanced course \pause 
\item Seminar work and presentations by \LaTeX  \pause 
\item Beamerclass
\end{enumerate}
\end{frame}

\section{Section Nr.3} 
\subsection{Tables}
\begin{frame}
\frametitle{Table}
\begin{tabular}{|c|c|c|}
\hline
\textbf{Date} & \textbf{Presentator} & \textbf{Title} \\
\hline
WS 04/05 & C.S. &  First Steps with \LaTeX  \\
\hline
SS 05 & C.S. & \LaTeX \ Course continues \\
\hline
\end{tabular}
\end{frame}


\begin{frame}
\frametitle{Tables with pauses}
\begin{tabular}{c c c}
A & B & C \\ 
\pause 
1 & 2 & 3 \\  
\pause 
A & B & C \\ 
\end{tabular} 
\end{frame}


\section{Section Nr.4}
\subsection{Blocks}
\begin{frame}\frametitle{Blocks}

\begin{block}{Block title}
Block text
\end{block}

\begin{exampleblock}{Block title}
Block text 
\end{exampleblock}


\begin{alertblock}{Block title}
Block text 
\end{alertblock}
\end{frame}


\subsection{The $\backslash$Only command}
\begin{frame}[fragile]\frametitle{Elements of fixed position}
\begin{verbatim}...\usepackage[absolute,overlay]{textpos} \end{verbatim}
\vspace*{3cm}
\begin{textblock}{0.6}(0.1,0.5)
\only<1>{\[A=B+C \] \[B=C+D\]}
\only<2>{\[A=2C+D\]}
\end{textblock}
\end{frame}


\section[Bibliography]{Literature}
\begin{frame}\frametitle{Sources \& Literature}

\begin{thebibliography}{9}
\bibitem[Beamerpaket]{paket} \emph{Beamer Paket} \\ 
\text{http://latex-beamer.sourceforge.net/}
\bibitem[Beamerdokumentation]{doku} \emph{User's Guide to the Beamer} 
\bibitem[Dante]{dante} \emph{DANTE e.V.} \text{http://www.dante.de}   
\end{thebibliography}


\end{frame}



\end{document}

